\begingroup
\shorthandoff{"<>"}
\setlength{\arrayrulewidth}{1pt}
\arrayrulecolor{vlepsbased}
\begin{longtable}{|p{15cm}|}
  \hline
  \rowcolor{maincolor} 
  \multicolumn{1}{|l|}{\textcolor{white}{\textbf{Caso de Uso}}} \\
  \hline
  \endfirsthead
  \multicolumn{1}{c}{{\tablename\ \thetable{} -- Continuación}} \\
  \hline
  \rowcolor{maincolor} 
  \multicolumn{1}{|l|}{\textcolor{white}{\textbf{Caso de Uso (Continuación)}}} \\
  \hline
  \endhead
  \hline \multicolumn{1}{|r|}{\textit{Continúa en la siguiente página...}} \\ \hline
  \endfoot
  \hline
  \caption{Caso de Uso CU8 — Cambio de contraseña.}
  \label{TB:CU8_PERFIL_PASSWORD}
  \endlastfoot

  \rowcolor{ulepsbased}
  \begin{tabularx}{\linewidth}{@{}p{7.5cm}p{7.5cm}@{}}
    \textbf{Identificador:} CU8 & \textbf{Actor principal:} Usuario
  \end{tabularx} \\ \hline

  \textbf{Nombre:} Cambio de contraseña \\ \hline

  \rowcolor{ulepsbased}
  \textbf{Descripción:} \newline 
  El usuario decide actualizar su contraseña de acceso desde su perfil introduciendo la clave actual y su nueva contraseña. \\ \hline

  \textbf{Precondiciones:} \newline
  1. El usuario ha iniciado sesión y tiene un token válido. \\ \hline

  \rowcolor{ulepsbased}
  \textbf{Flujo principal:} \newline
  1. El usuario navega a la sección de "Perfil / Ajustes" y selecciona la opción para cambiar contraseña. \newline
  2. El usuario introduce su contraseña actual y una nueva contraseña. \newline
  3. El usuario pulsa el botón de confirmación ("Cambiar contraseña"). \newline
  4. La interfaz envía una solicitud POST con las contraseñas al backend. \newline
  5. El sistema valida los permisos (mediante el token JWT) y comprueba las políticas de seguridad de la nueva clave. \newline
  6. El sistema verifica que la contraseña actual proporcionada coincida con el hash almacenado en la base de datos. \newline
  7. El sistema genera un nuevo hash para la contraseña y lo actualiza en la base de datos. \newline
  8. El backend retorna un código 204 No Content confirmando la operación. \newline
  9. La interfaz muestra un aviso de éxito notificando el cambio de credenciales. \\ \hline

  \textbf{Flujos alternativos:} \newline
  \textbf{FA-01:} Contraseña actual incorrecta: \newline
  1. El sistema recupera el hash almacenado del usuario y comprueba que no coincide con la "contraseña actual" proporcionada. \newline
  2. Se rechaza la operación, retornando HTTP 400 Bad Request. \newline
  3. La interfaz informa al usuario de que sus credenciales son erróneas. \newline
  \textbf{FA-02:} Nueva contraseña inválida: \newline
  1. El backend detecta que la nueva contraseña no cumple con las políticas de seguridad mínimas (longitud, caracteres) o el frontend detecta que no coinciden (si hubiera confirmación). \newline
  2. Se aborta la operación con un error de validación (HTTP 400). \newline
  3. Se solicita al usuario que elija una contraseña diferente. \\ \hline

  \rowcolor{ulepsbased}
  \textbf{Postcondiciones:} \newline
  1. El hash de la contraseña del usuario se actualiza en el sistema y se requerirá la nueva clave en su próximo inicio de sesión. \\ \hline

\end{longtable}
\endgroup
